%% ============================================================
%% Figures — embedded at end of body, before bibliography.
%% Captions mirror the Markdown source; the \ref labels match the
%% sed-applied transforms in the Makefile that turn "Figure N" into
%% \ref{fig:...} in the body text.
%% ============================================================

\begin{figure}[!htbp]
\centering
\includegraphics[width=0.85\textwidth]{figures/fig01_scheme.pdf}
\caption{Well-balanced finite-volume scheme at an x-face between cells L and R over a stepped bed (to scale). The cells carry actual depths $h_L, h_R$ (blue brackets at the cell centres) with free surfaces $\eta _L, \eta _R$. The Audusse (2004) hydrostatic reconstruction defines $z_max = max(z_L, z_R)$ (purple dashed) and the reconstructed face depths $h*_L = max(\eta _L - z_max, 0)$, $h*_R = max(\eta _R - z_max, 0)$ (purple brackets), on which the HLLC flux $F$ (vermillion) is evaluated. The bed is reconstructed linearly to the shared face value $z_face = ½(z_L + z_R)$, which makes the Audusse pressure correction vanish and moves the bed-slope force into the cell-centred source term (orange) --- the construction that gives the exact lake-at-rest balance verified in \S3.1. Schematic (no simulation data); generated by \texttt{fig01\_scheme.R}.}
\label{fig:scheme}
\end{figure}

\begin{figure}[!htbp]
\centering
\includegraphics[width=\textwidth]{figures/fig02_verification.pdf}
\caption{Verification against analytical references: simulated profiles (points) overlaid on the analytical solution (line). (a) Stoker/Ritter dam-break on a dry bed at $t = 4$ s --- the rarefaction fan and the leading dry front are captured (400-cell channel, mid-row slice). (b) MacDonald uniform flow at steady state, y- axis zoomed to ±1.5 \% of the normal depth \texttt{h\_n} to expose the well-balanced preservation (the profile stays flat at \texttt{h\_n} on a sloped bed under Manning friction). (c) Thacker oscillating paraboloid at $t = T/2$, centre-row slice --- the parabolic cap and its moving wet/dry shoreline track the analytical solution on a curved bed. Quantitative error metrics for all benchmarks are in Table 1 (computed over the benchmark-specific regions: the Stoker rarefaction fan, the 2D Thacker interior); they are not annotated here to avoid conflating the 1D-slice view of this figure with the table's domain-wide norms. Generated by \texttt{fig02\_verification.R} from the \texttt{gen\_verification\_data} example output.}
\label{fig:verification}
\end{figure}

\begin{figure}[!htbp]
\centering
\includegraphics[width=\textwidth]{figures/fig03_uk_ea_t6.pdf}
\caption{UK Environment Agency 2D benchmark Test 6: urban dam-break depth field at $t = 30$ s on the 500 $\times$ 100 m domain (250 $\times$ 50 cells, $\Delta x = \Delta y = 2$ m). The reservoir ($x < 100$ m, dashed dam line) collapses into the dry downstream and reaches the first building cluster ($x \approx 150$--170 m) by $t = 30$ s, wrapping around the obstacles with the expected wakes and inter-building jets. The six raised-bed buildings (grey) remain dry throughout --- water navigates the array without overtopping them --- while the wet/dry front is captured without spurious oscillation. Depth on a perceptually uniform sequential colour scale (scico *devon*); downstream buildings ($x \geq 250$ m) are still dry at this time as the front has not yet arrived. Generated by \texttt{fig03\_uk\_ea\_t6.R} from the \texttt{gen\_verification\_data} example output.}
\label{fig:ukea_t6}
\end{figure}

\begin{figure}[!htbp]
\centering
\includegraphics[width=\textwidth]{figures/fig04_huasco_application.pdf}
\caption{Río Huasco reach under the regulated 2017 release, 200 $\times$ 67-cell 30 m subset (UTM 19S), at the day-11 peak. Five panels sharing the reach extent: (a) ESA WorldCover 2021 land cover --- riparian tree and shrub vegetation tracks the thalweg within bare Atacama hillslopes; (b) the derived Manning field $n(x, y)$ mapping land cover to roughness ($n = 0.025$ bare to $0.10$ tree); (c) inundation depth with a single uniform $n = 0.04$; (d) inundation depth with the variable $n(x, y)$; (e) the difference $\Delta h = (d) - (c)$, hill-shaded base, divergent scale. Panels (c)--(e) share a hillshade underlay. The positive $\Delta h$ (warm) concentrated in the channel shows the riparian roughness deepening and retaining the flow --- the +0.33 m mean channel deepening and +29 \% retained volume reported in \S4.3. Generated by \texttt{fig04\_huasco\_application.R}, which reuses the solver-2d example rasters (\texttt{huasco\_subset\_{dem,landcover}.tif}, \texttt{huasco\_2d\_depth\_day\_11{,\_landcover}.tif}).}
\label{fig:huasco}
\end{figure}

\begin{figure}[!htbp]
\centering
\includegraphics[width=0.75\textwidth]{figures/fig05_convergence.pdf}
\caption{Mesh-refinement convergence on the Thacker oscillation (\S3.7). (a) Relative L1 (circles) and L2 (triangles) error in depth at $t = T/2$ versus cell size $\Delta x$, log-log, with order-1 and order-2 reference slopes (dashed); the overall fitted slopes are 1.73 (L1) and 1.58 (L2). (b) Order observed between successive meshes, $\log_2(e_{\text{coarse}}/e_{\text{fine}})$, plotted against the finer mesh of each pair. Panel (b) carries the claim panel (a) cannot: a single fitted slope averages the degradation away, whereas the pairwise orders show it directly, falling from 1.81 to 1.59 in L1 and from 1.71 to 1.43 in L2 as the moving wet/dry shoreline --- locally first order --- comes to dominate the shrinking smooth-region error. Generated by \texttt{fig05\_convergence.R} from the \texttt{gen\_convergence} example.}
\label{fig:convergence}
\end{figure}

\begin{figure}[!htbp]
\centering
\includegraphics[width=0.9\textwidth]{figures/fig06_head_to_head.pdf}
\caption{Head-to-head against ANUGA \citep{Roberts2015} on the Stoker/Ritter dam-break at matched resolution ($\Delta x = 1$ m, \S3.8). Both solvers (circles: hydroflux; triangles: ANUGA) overlaid on the analytical Ritter profile (line) at $t = 4$ s; the rarefaction fan $[37.5, 75.1]$ m is shaded. The dashed vertical marks the initial dam at $x = 50$ m. Relative error norms over the fan, annotated: L1 4.1 \% / L2 3.6 \% / $L^{\infty}$ 5.3 \% $h_L$ for hydroflux versus 2.6 \% / 2.7 \% / 4.4 \% for ANUGA. Both stay below 6 \% at this coarse resolution; hydroflux's reference 400-cell run (Table 1) reaches L1 1.0 \%. Generated by \texttt{fig06\_head\_to\_head.R} from the \texttt{anuga\_stoker\_compare.py} + \texttt{gen\_stoker\_coarse} example outputs.}
\label{fig:head_to_head}
\end{figure}
